\begin{proposition}[Counterexample: The universal two-vector-field bound]\hypertarget{tag-01SX}{}\label{inv:E2.4}
\StatementRecord{inv:E2.4}{proved}{01SX}{\SeriesRef{inv:E2.2}; \citet[Corollary~1.3, preprint p.~2]{HKP2000}; \citet[Remarks 1.2 and 9.26]{AC2026}}

The assertion that every complex structure on $S^6$ has
$h^0(T)\le2$, printed in
\citet[Corollary~1.3, preprint p.~2]{HKP2000}
and repeated in \citet[Remarks 1.2 and 9.26]{AC2026}, is false
for the comb members.
\end{proposition}


\begin{proof}
Proposition~\ref{inv:E2.2} gives $h^0(T_{X_m})=2m+1\ge25$ for $m\ge12$.
In the proof of \citet[Corollary~1.3, preprint p.~2]{HKP2000},
a third field is expressed as a meromorphic linear combination of two fields;
the coefficients are then declared constant using the
Campana--Demailly--Peternell assertion that every meromorphic function
on a complex six-sphere is constant. For the combs, the surjective algebraic reduction
$f_m$ induces an injection $f_m^*:\C(t)\hookrightarrow\mathcal M(X_m)$,
so its pullback of $t$ is nonconstant.
\end{proof}